\begin{thebibliography}{99}

\bibitem{ref:tdf1555} 国务院. 关于印发《“十五五”碳达峰行动方案》的通知[EB/OL]. 国发〔2026〕22号. (2026-07-09)[2026-09-09].
\url{https://www.gov.cn/zhengce/content/202607/content_7074826.htm}.

\bibitem{ref:wlw2026} 国家发展改革委, 交通运输部. 关于印发《物流网建设实施方案》的通知[EB/OL]. 发改经贸〔2026〕1241号. (2026-08-18)[2026-09-09].
\url{https://zfxxgk.ndrc.gov.cn/web/iteminfo.jsp?id=20659}.

\bibitem{ref:6} 李得成, 陈彦如, 张宗成. 基于分支定价算法的电动车与燃油车混合车辆路径问题研究[J]. 系统工程理论与实践, 2021, 41(4): 995--1009.\newline
Li D C, Chen Y R, Zhang Z C. A branch-and-price algorithm for electric vehicle routing problem with time windows and mixed fleet[J]. Systems Engineering --- Theory \& Practice, 2021, 41(4): 995--1009.

\bibitem{ref:23} 陈婉茹, 徐光明, 张得志, 等. 碳交易机制下多中心混合车队配送路径和速度优化研究[J]. 系统工程理论与实践, 2023, 43(11): 3320--3335.\newline
Chen W R, Xu G M, Zhang D Z, et al. Multi-depot mixed fleet routing and speed optimization under a carbon trading mechanism[J]. Systems Engineering --- Theory \& Practice, 2023, 43(11): 3320--3335.

\bibitem{ref:qiu2024} Qiu Y, Ding S, Pardalos P M. Routing a mixed fleet of electric and conventional vehicles under regulations of carbon emissions[J]. International Journal of Production Research, 2024, 62(16): 5720--5736.

\bibitem{ref:bruglieri2025} Bruglieri M, \c{C}atay B, Keskin M, et al. The mixed-fleet vehicle routing problem with low emission zones[J]. Transportation Research Part E: Logistics and Transportation Review, 2025, 201: 104230.

\bibitem{ref:gao2025} 高咏玲, 乔源, 徐猛. 考虑电动车配置与路径优化的城市配送车队更新政策效应研究[J]. 中国管理科学, 2025, 33(8): 156--165.\newline
Gao Y L, Qiao Y, Xu M. Effects analyses of delivery fleet renewal policies considering electric vehicle deployment and routing optimization[J]. Chinese Journal of Management Science, 2025, 33(8): 156--165.

\bibitem{ref:dantzig} Dantzig G B, Ramser J H. The truck dispatching problem[J]. Management Science, 1959, 6(1): 80--91.

\bibitem{ref:27} Goeke D, Schneider M. Routing a mixed fleet of electric and conventional vehicles[J]. European Journal of Operational Research, 2015, 245(1): 81--99.

\bibitem{ref:45} Keskin M, \c{C}atay B. Partial recharge strategies for the electric vehicle routing problem with time windows[J]. Transportation Research Part C: Emerging Technologies, 2016, 65: 111--127.

\bibitem{ref:montoya} Montoya A, Guéret C, Mendoza J E, et al. The electric vehicle routing problem with nonlinear charging function[J]. Transportation Research Part B: Methodological, 2017, 103: 87--110.

\bibitem{ref:12} Froger A, Jabali O, Mendoza J E, et al. The electric vehicle routing problem with capacitated charging stations[J]. Transportation Science, 2022, 56(2): 460--482.

\bibitem{ref:zhou2026} 周鲜成, 李松明, 王莉, 等. 考虑非线性能耗的时间依赖型电动车辆路径模型及改进的鲸鱼优化算法[J]. 中国管理科学, 2026, 34(3): 214--225.\newline
Zhou X C, Li S M, Wang L, et al. Research on time dependent electric vehicle routing model with nonlinear energy consumption and an improved whale optimization algorithm[J]. Chinese Journal of Management Science, 2026, 34(3): 214--225.

\bibitem{ref:deng} Deng J, Hu H, Gong S, et al. Impacts of charging pricing schemes on cost-optimal logistics electric vehicle fleet operation[J]. Transportation Research Part D: Transport and Environment, 2022, 109: 103333.

\bibitem{ref:shi2025} Shi W, Wang N, Zhou L, et al. The bi-objective mixed-fleet vehicle routing problem under decentralized collaboration and time-of-use prices[J]. Expert Systems with Applications, 2025, 273: 126875.

\bibitem{ref:52} Li Z, Chen Z, Li H, et al. On the value of orderly electric vehicle charging in carbon emission reduction[J]. Transportation Research Part D: Transport and Environment, 2024, 135: 104383.

\bibitem{ref:du2025} Du B, Jia H, Zhang B, et al. Evaluating the carbon-emission reduction potential of employing low-carbon demand response to guide electric-vehicle charging: A Chinese case study[J]. Applied Energy, 2025, 397: 126196.

\bibitem{ref:martin2025} Martin S, Powell S, Rajagopal R. Cascading marginal emissions signals for green charging with growing electric vehicle adoption[J]. Nature Communications, 2025, 16: 10150.

\bibitem{ref:44} Zhen L, Ma C, Wang K, et al. Multi-depot multi-trip vehicle routing problem with time windows and release dates[J]. Transportation Research Part E: Logistics and Transportation Review, 2020, 135: 101866.

\bibitem{ref:sahin} Şahin M K, Yaman H. A branch and price algorithm for the heterogeneous fleet multi-depot multi-trip vehicle routing problem with time windows[J]. Transportation Science, 2022, 56(6): 1636--1657.

\bibitem{ref:zhao2024} Zhao J, Poon M, Tan V Y F, et al. A hybrid genetic search and dynamic programming-based split algorithm for the multi-trip time-dependent vehicle routing problem[J]. European Journal of Operational Research, 2024, 317(3): 921--935.

\bibitem{ref:49} Fern\'{a}ndez E, Roca-Riu M, Speranza M G. The shared customer collaboration vehicle routing problem[J]. European Journal of Operational Research, 2018, 265(3): 1078--1093.

\bibitem{ref:7} Wang Y, Zhou J, Sun Y, et al. Collaborative multidepot electric vehicle routing problem with time windows and shared charging stations[J]. Expert Systems with Applications, 2023, 219: 119654.

\bibitem{ref:wang2024} Wang Y, Wei Z, Luo S, et al. Collaboration and resource sharing in the multidepot time-dependent vehicle routing problem with time windows[J]. Transportation Research Part E: Logistics and Transportation Review, 2024, 192: 103798.

\bibitem{ref:8} Soriano A, Gansterer M, Hartl R F. The multi-depot vehicle routing problem with profit fairness[J]. International Journal of Production Economics, 2023, 255: 108669.

\bibitem{ref:91} 饶卫振, 苗晓河, 朱庆华, 等. 考虑企业服务质量差异的协作配送问题及成本分摊方法研究[J]. 系统工程理论与实践, 2022, 42(10): 2721--2739.\newline
Rao W Z, Miao X H, Zhu Q H, et al. Research on collaborative distribution problems and cost allocation methods considering differences in service quality of logistics enterprises[J]. Systems Engineering --- Theory \& Practice, 2022, 42(10): 2721--2739.

\bibitem{ref:14} Psaraftis H N. Dynamic vehicle routing problems[M]//Golden B L, Assad A A. Vehicle Routing: Methods and Studies. Amsterdam: North-Holland, 1988: 223--248.

\bibitem{ref:18} 李阳, 范厚明, 张晓楠. 动态需求下车辆路径问题的周期性优化模型及求解[J]. 中国管理科学, 2022, 30(8): 254--266.\newline
Li Y, Fan H M, Zhang X N. A periodic optimization model and solution for capacitated vehicle routing problem with dynamic requests[J]. Chinese Journal of Management Science, 2022, 30(8): 254--266.

\bibitem{ref:51} 姜广田, 纪皎月, 董佳伟. 绿色物流配送下的多车型动态车辆路径优化[J]. 系统工程理论与实践, 2024, 44(7): 2362--2380.\newline
Jiang G T, Ji J Y, Dong J W. Multi-vehicle dynamic vehicle routing optimization in green logistics distribution[J]. Systems Engineering --- Theory \& Practice, 2024, 44(7): 2362--2380.

\bibitem{ref:71} 邱莹莹. 低碳背景下考虑动态需求的混合车队物流配送路径问题研究[D]. 上海: 上海理工大学, 2024.\newline
Qiu Y Y. Research on the vehicle routing problem of mixed-fleet considering dynamic demand under the background of low carbon[D]. Shanghai: University of Shanghai for Science and Technology, 2024.

\bibitem{ref:mardesic2024} Marde\v{s}i\'{c} N, Erdeli\'{c} T, Cari\'{c} T, et al. Review of stochastic dynamic vehicle routing in the evolving urban logistics environment[J]. Mathematics, 2024, 12(1): 28.

\bibitem{ref:shahbazian2025} Shahbazian R, Ciacco A, Macrina G, et al. Knowledge-guided hybrid deep reinforcement learning for the dynamic multi-depot electric vehicle routing problem[J]. Computers \& Operations Research, 2025, 184: 107217.

\bibitem{ref:94} Wang W, Adulyasak Y, Cordeau J F. Electric vehicle routing with heterogeneous charging stations[J]. Computers \& Operations Research, 2026, 189: 107374.

\bibitem{ref:95} 巩亮, 郑世龙, 许得杰, 等. 低碳视角下油电混合车队配送路径优化研究[J]. 控制与决策, 2026, 41(5): 1338--1347.\newline
Gong L, Zheng S L, Xu D J, et al. Optimization of distribution paths for hybrid fuel-electric fleet under low-carbon perspective[J]. Control and Decision, 2026, 41(5): 1338--1347.

\bibitem{ref:76} Vidal T, Crainic T G, Gendreau M, et al. A hybrid genetic algorithm for multidepot and periodic vehicle routing problems[J]. Operations Research, 2012, 60(3): 611--624.

\bibitem{ref:77} Vidal T, Crainic T G, Gendreau M, et al. A unified solution framework for multi-attribute vehicle routing problems[J]. European Journal of Operational Research, 2014, 234(3): 658--673.

\bibitem{ref:hgs-cvrp-2022} Vidal T. Hybrid genetic search for the CVRP: Open-source implementation and SWAP* neighborhood[J]. Computers \& Operations Research, 2022, 140: 105643.

\bibitem{ref:64} Wouda N A, Lan L, Kool W. PyVRP: A high-performance VRP solver package[J]. INFORMS Journal on Computing, 2024, 36(4): 943--955.

\bibitem{ref:changjiang2024} 邬博华, 曹海花, 叶之楠. 经济性驱动，新能源轻重卡兑现爆发式增长[R]. 武汉: 长江证券研究所, 2024: 5.

\bibitem{ref:bj-tou} 北京市发展和改革委员会. 关于进一步完善本市分时电价机制等有关事项的通知[EB/OL]. 京发改规〔2023〕11号. (2023-08-21)[2026-09-09].
\url{https://fgw.beijing.gov.cn/fgwzwgk/2024zcwj/bwgfxwj/202308/t20230821_3718725.htm}.

\bibitem{ref:bj-fee} 北京市发展和改革委员会. 关于本市电动汽车充电服务收费有关问题的通知[EB/OL]. 京发改〔2015〕848号. (2015-04-24)[2026-09-09].
\url{https://www.beijing.gov.cn/zhengce/zhengcefagui/201905/t20190522_58449.html}.

\bibitem{ref:mee2025} 生态环境部. 全国碳市场发展报告（2025）[R]. 北京: 生态环境部, 2025.

\bibitem{ref:zhang2025} Zhang J. Pickup and delivery planning for the crowdsourced freight delivery routing problem[J]. PLoS ONE, 2025, 20(2): e0318432.

\bibitem{ref:osrm} Luxen D, Vetter C. Real-time routing with OpenStreetMap data[C]//Proceedings of the 19th ACM SIGSPATIAL International Conference on Advances in Geographic Information Systems. New York: ACM, 2011: 513--516.

\bibitem{ref:82} Li Y, Zhang S, Li W, et al. High temporal and spatial resolution projected electricity carbon emission factors of China from 2025--2060[J]. Scientific Data, 2026, 13: 926.

\bibitem{ref:75} Lei Z, Hao J K. A GPU-accelerated hybrid method for a class of multi-depot vehicle routing problems[EB/OL]. arXiv:2605.05208. (2026-02-24)[2026-09-09].
\url{https://arxiv.org/abs/2605.05208}.

\bibitem{ref:74} Vidal T, Crainic T G, Gendreau M, et al. A hybrid genetic algorithm with adaptive diversity management for a large class of vehicle routing problems with time-windows[J]. Computers \& Operations Research, 2013, 40(1): 475--489.

\bibitem{ref:woody2021} Woody M, Vaishnav P, Craig M T, et al. Charging strategies to minimize greenhouse gas emissions of electrified delivery vehicles[J]. Environmental Science \& Technology, 2021, 55(14): 10108--10120.

\bibitem{ref:ndrc-tou} 国家发展改革委. 关于进一步完善分时电价机制的通知[EB/OL]. 发改价格〔2021〕1093号. (2021-07-26)[2026-09-09].
\url{https://www.ndrc.gov.cn/xxgk/zcfb/tz/202107/t20210729_1292067.html}.

\bibitem{ref:hebei-tou} 河北省发展和改革委员会. 关于进一步完善河北南网工商业及其他用户分时电价政策的通知[EB/OL]. 冀发改能价〔2022〕1364号. (2022-10-28)[2026-09-09].
\url{http://info.hebei.gov.cn/hbszfxxgk/329975/329988/330035/6852718/7041023/index.html}.

\end{thebibliography}
